\section{Information, terminal width, and risk}
\label{sec:risk_setup}
Let $F\subset\mathbb R^n$ be nonempty, compact, and convex, and let
$T(y)=c^\top y+d$ be the scalar quantity to be audited. A fixed finite
collection of exact linear measurements is available. Write $L_Q$ for
the measurement matrix of a selected set $Q$, and let $\mathcal Q$ be
the nonempty finite family of admissible final sets, available to both
policy classes. Admissibility is determined before observing answers; for example,
$\sum_{j\in Q}k_j\le B$ for nonnegative fixed costs $k_j$.
Coordinate observations correspond to rows of the identity matrix.
The results do not require nonnegative targets or disjoint aggregates.

For a feasible record $y$, define the information fiber and its sharp
target width by
\begin{equation}
 F_Q(y)=\{z\in F:L_Qz=L_Qy\},\qquad
 W_Q(y)=\max_{z\in F_Q(y)}T(z)-\min_{z\in F_Q(y)}T(z).
 \label{eq:risk_fiber}
\end{equation}
The extrema exist by compactness, and $0\le W_Q(y)\le D_T$, where
$D_T=\max_F T-\min_F T$. This is the local diameter of the
scalar information fiber. The optimized quantity is the full fiber's
diameter, rather than the actual error of a selected point estimate;
it is not a confidence interval. An additive constant in $T$ does
not affect $W_Q$. In particular, a paired squared-loss contrast is
affine in the unknown verification targets and fits this formulation.

A randomized batch policy draws $Q\sim\pi$ before seeing answers. An
adaptive policy selects successive measurements and a stopping decision
from the answers already observed and an independent internal random
seed. Every execution must terminate after finitely many queries at a
set in $\mathcal Q$. Query and stopping rules are jointly measurable in
the observed history and seed; observations are exact. For a finite-valued,
nondecreasing function $\phi:[0,\infty)\to\mathbb R$, define
\begin{equation}
 R_{\mathrm{batch}}^{\phi}
 =\inf_{\pi\in\Delta(\mathcal Q)}\sup_{y\in F}
   \mathbb E_{Q\sim\pi}\phi(W_Q(y)),\qquad
 R_{\mathrm{adaptive}}^{\phi}
 =\inf_{\mathcal P}\sup_{y\in F}
   \mathbb E\phi(W_{Q_{\mathcal P}(y)}(y)).
 \label{eq:risk_objectives}
\end{equation}
The hidden record is fixed before the random seed is drawn. Thus the
supremum is outside the expectation. When measurement error is represented
by additional hidden coordinates, one fixed error vector is part of that
record for the entire execution. An adversary choosing a new error after
each query is a different information model. These objectives concern terminal
width under a hard budget, rather than expected requests until a winner
is certified. The expectations are well defined: the widths are Borel
measurable by Proposition~\ref{prop:risk_concavity_transcript}, and a
finite-valued monotone $\phi$ is Borel measurable and bounded on $[0,D_T]$.

For a required precision $\tau\geq0$, let $\mathcal C$ be either the
randomized batch or randomized adaptive policy class, and write
$Q_{\mathcal P}(y,\omega)$ for the terminal set under internal seed
$\omega$. Define its optimal worst-case precision-success probability by
\begin{equation}
\begin{aligned}
 S_\tau^{\mathcal C}
 &=\sup_{\mathcal P\in\mathcal C}\inf_{y\in F}
   \Pr_\omega\!\left\{W_{Q_{\mathcal P}(y,\omega)}(y)\leq\tau\right\}\\
 &=1-R_{\mathcal C}^{\phi_\tau},\qquad
   \phi_\tau(w)=\boldsymbol{1}\{w>\tau\},
\end{aligned}
\label{eq:risk_precision_success}
\end{equation}
where $R_{\mathcal C}^{\phi_\tau}$ denotes the corresponding risk in
\eqref{eq:risk_objectives}. The identity follows by complementing the
success event and reversing each supremum or infimum under subtraction
from one; it does not require an optimal policy to attain either bound.
The probability concerns the algorithm's seed at a fixed hidden record,
and equality at the threshold counts as success. Reaching this precision
does not require the interval to exclude zero, so it need not certify a
winning model.

\begin{proposition}[Concavity and the transcript fiber]
\label{prop:risk_concavity_transcript}
For every fixed $Q$, $W_Q$ is nonnegative, finite, concave, and upper
semicontinuous on $F$. For every adaptive execution, its terminal
information fiber is precisely $F_Q(y)$ for its terminal queried set $Q$.
\end{proposition}
\begin{proof}
Choose maximizing and minimizing endpoints $u_k,v_k\in F_Q(y_k)$,
$k=1,2$. For $0\le\lambda\le1$, convexity of $F$ and linearity of the
measurements make $\lambda u_1+(1-\lambda)u_2$ and
$\lambda v_1+(1-\lambda)v_2$ feasible endpoints at
$\lambda y_1+(1-\lambda)y_2$. Their target difference is
$\lambda W_Q(y_1)+(1-\lambda)W_Q(y_2)$, proving concavity.
For upper semicontinuity, take $y_k\to y$ and a subsequence attaining
the limsup of the widths. Compactness supplies a further subsequence
whose endpoint pairs converge to $u,v\in F_Q(y)$. Its limiting
difference is at most $W_Q(y)$.

For the transcript claim, fix the internal random seed. If no query is
made, both fibers are $F$. Otherwise, any
$z\in F_Q(y)$ gives the same answer to the first queried measurement.
Inductively, equal previous answers produce the same next measurement
and the same next answer. The stopping decision is also unchanged.
Thus every completion agreeing with all queried answers follows the
same entire transcript. Conversely, any completion consistent with
that transcript must agree on every queried answer. Query choices and
the stopping time therefore impose no additional restrictions on the
fiber beyond those exact observations. This argument is pointwise in the
seed; no conditioning on a positive-probability seed event is required.
\end{proof}

\section{Which width penalties can benefit from adaptation?}
\label{sec:risk_classification}
The next result uses classical minimax theory. The contribution is the
resulting classification for sharp information-fiber widths and its
coordinate-query converse, rather than a new minimax or duality theorem.
The proof has two parts. For concave penalties, a single hardest record
limits every final query set, including those chosen adaptively; Jensen's
inequality supplies a batch policy meeting that limit. For any nonconcave
penalty, a three-coordinate construction realizes a strict midpoint
violation as an adaptive advantage. These arguments establish the universal
claim; the later numerical checks test implementations on particular inputs.

\begin{proposition}[Universal classification of width penalties]
\label{prop:risk_classification}
Let $\phi:[0,\infty)\to\mathbb R$ be finite-valued and nondecreasing.
Randomized batch and randomized adaptive querying have equal minimax
$\phi$-risk for every problem satisfying Section~\ref{sec:risk_setup}
if and only if $\phi$ is concave. More precisely, put
\begin{equation}
 \rho=\min_{\pi\in\Delta(\mathcal Q)}\max_{y\in F}
         \sum_{Q\in\mathcal Q}\pi_Q W_Q(y).
 \label{eq:risk_rho}
\end{equation}
If $\phi$ is concave, then
\begin{equation}
 R_{\mathrm{adaptive}}^{\phi}
 =R_{\mathrm{batch}}^{\phi}=\phi(\rho).
 \label{eq:risk_concave_value}
\end{equation}
Every batch policy attaining~\eqref{eq:risk_rho} simultaneously attains
this equality for every such concave $\phi$. If $\phi$ is not concave,
a strict adaptive improvement already occurs on a three-coordinate
compact convex polytope, with all coordinates queryable at unit cost
and a hard budget of two.
\end{proposition}
\begin{proof}
The function $\sum_Q\pi_QW_Q(y)$ is affine and continuous in $\pi$,
and concave and upper semicontinuous in $y$. Sion's minimax theorem
therefore yields
\begin{equation}
 \rho=\max_{y\in F}\min_{Q\in\mathcal Q}W_Q(y).
 \label{eq:risk_hardest_record}
\end{equation}
All values are finite. Writing $G(\pi)=\max_{y\in F}\sum_Q\pi_QW_Q(y)$,
the bound $|G(\pi)-G(\pi')|\le D_T\|\pi-\pi'\|_1$ proves continuity.
Thus the minimum in~\eqref{eq:risk_rho} is attained on the compact simplex.
The maximum in~\eqref{eq:risk_hardest_record} is attained because the
minimum of finitely many upper semicontinuous functions is upper
semicontinuous on compact $F$. Hence there is one record $y^\star\in F$
such that $W_Q(y^\star)\ge\rho$ for every admissible $Q$.
By Proposition~\ref{prop:risk_concavity_transcript}, every execution of
every adaptive policy has terminal width at least $\rho$ at this same
record. Monotonicity gives the lower bound $\phi(\rho)$.

Let $\pi^\star$ attain~\eqref{eq:risk_rho}. Concavity, Jensen's
inequality, and monotonicity give, for every $y\in F$,
\[
 \sum_Q\pi^\star_Q\phi(W_Q(y))
 \le \phi\!\left(\sum_Q\pi^\star_QW_Q(y)\right)
 \le\phi(\rho).
\]
This proves the upper bound and simultaneous optimality. No continuity
of $\phi$ at zero is needed; the finite-mixture form of Jensen's
inequality suffices. For example, $\phi(w)=\boldsymbol{1}\{w>0\}$ is
concave on $[0,\infty)$ despite its jump at zero.

For necessity, a nonconcave, nondecreasing finite-valued $\phi$ must
have a strict midpoint-concavity violation: there exist $0\le l<h$ such
that, with $m=(l+h)/2$,
\begin{equation}
 \phi(m)<\{\phi(l)+\phi(h)\}/2.
 \label{eq:risk_midpoint_violation}
\end{equation}
Indeed, if midpoint concavity held, for $0\le x\le z$ and dyadic
$q_k\uparrow\theta$ its dyadic inequalities and monotonicity would give
\[
\phi((1-\theta)x+\theta z)
\ge\phi((1-q_k)x+q_kz)
\ge(1-q_k)\phi(x)+q_k\phi(z).
\]
Taking the limit only on the right proves full concavity without a
continuity assumption, a contradiction. Set $d_0=h-l$ and consider
\begin{equation}
 F_{l,h}=\{(a,b,c):0\le a\le1,\ 0\le b\le l+d_0a,
                  \ 0\le c\le h-d_0a\},\qquad T=ha+b+c.
 \label{eq:risk_converse_polytope}
\end{equation}
All three coordinates are available at unit cost, with budget two.
The widths after querying $\{a,b\}$ and $\{a,c\}$ are respectively
$h-d_0a$ and $l+d_0a$. An adaptive policy first observes $a$ and then
queries the coordinate with the larger remaining capacity. Its terminal
width is at most $m$, so its risk is at most $\phi(m)$.
At $(a,b,c)=(1/2,0,0)$, these two widths equal $m$, while the width
after querying $\{b,c\}$ is $h$. Thus every admissible terminal set
has width at least $m$, and the adaptive optimum is exactly $\phi(m)$.

Now put equal adversarial mass on $(0,0,0)$ and $(1,0,0)$.
Each batch $\{a,b\}$ or $\{a,c\}$ has average penalty
$\{\phi(l)+\phi(h)\}/2$. The batch $\{b,c\}$ has width $h$ at both
records, so its average penalty is no smaller. A batch using fewer
coordinates cannot improve upon an extension to two coordinates,
because more observations only shrink the fiber. Every randomized
batch policy consequently has worst-case risk at least
$\{\phi(l)+\phi(h)\}/2>\phi(m)$. This is the required strict separation.
\end{proof}

The universal quantifier is essential. A nonconcave penalty permits an
adaptive advantage somewhere in this class; it need not yield an
advantage for every feasible set. The theorem covers heterogeneous fixed
costs and overlapping constraints at the level of information theory.
It does not imply a polynomial-time optimizer for an arbitrary feasible
set or budget.

\subsection{A three-coordinate example}
\label{sec:risk_small_example}
Take $0\le a\le1$, $0\le b\le a$, $0\le c\le1-a$, and
$T=a+b+c$, with two unit-cost queries. The widths for batches
$\{a,b\}$, $\{a,c\}$, and $\{b,c\}$ are $1-a$, $a$, and $1-b-c$.
Querying $a$ first and then $c$ if $a\le1/2$, or $b$ otherwise, leaves
width $\min(a,1-a)\le1/2$.

\begin{table}[htbp]
\centering
\caption{Exact constructed-example values. These are mathematical
guarantees, not forecasting improvements measured on public data.}
\label{tab:risk_three_coordinate}
\begin{tabular}{lcc}
\toprule
Criterion & Best randomized batch & Best adaptive policy\\
\midrule
Worst expected width & $1/2$ & $1/2$\\
Worst expected squared width & $1/2$ & $1/4$\\
Worst-case success, width $\le3/4$ & $1/2$ & $1$\\
\bottomrule
\end{tabular}
\end{table}

An equal mixture of $\{a,b\}$ and $\{a,c\}$ attains the batch values.
The endpoint records $(0,0,0)$ and $(1,0,0)$ give the corresponding
batch lower bounds. At $(1/2,0,0)$ every admissible terminal set has
width at least $1/2$, establishing the adaptive moment lower bounds.
The precision separation holds for every threshold $1/2\le\tau<1$.
Reaching a width threshold is distinct from certifying the sign of $T$:
a narrow sharp interval may still contain zero.

\section{Growing moment and precision separations}
\label{sec:risk_growing}
For $K\ge2$, let $a_1,\ldots,a_K\ge0$ sum to one and impose
$0\le b_j\le a_j$. The actual queryable coordinates are
$a_1,\ldots,a_{K-1},b_1,\ldots,b_K$; $a_K$ is determined by the
first $K-1$ capacities. Define
\begin{equation}
 T=\sum_{i<K}a_i+\sum_{j=1}^K b_j,\qquad B=2K-2.
 \label{eq:risk_K_setup}
\end{equation}
This is a compact convex polytope in $2K-1$ dimensions, and every
coordinate has the same unit cost.

\begin{proposition}[Sharp growing separations]
\label{prop:risk_K_family}
For~\eqref{eq:risk_K_setup}, the optimal worst expected width is $1/K$
for both randomized batch and adaptive policies. For every $p>1$,
\begin{equation}
 R_{\mathrm{batch}}^{w^p}=1/K,\qquad
 R_{\mathrm{adaptive}}^{w^p}=1/K^p.
 \label{eq:risk_K_moments}
\end{equation}
For required width $\tau=1/K$, the best worst-case success probabilities
are $1/K$ for randomized batches and $1$ for adaptive policies.
\end{proposition}
\begin{proof}
It suffices to consider batches using the full budget: adding observations
cannot increase width. Such a batch omits exactly one coordinate.
Omitting $b_j$ leaves width $a_j$; omitting $a_i$, $i<K$, leaves width
$a_i+a_K-b_i-b_K$. The latter follows because the remaining feasible
range of $a_i$ is $[b_i,a_i+a_K-b_K]$ and its coefficient in $T$ is one.
An adaptive policy queries all explicit capacities and then all $b_j$
except one having the smallest capacity. Its width is
$\min_j a_j\le1/K$. At the record with all $a_j=1/K$ and all $b_j=0$,
every full-budget batch has width at least $1/K$. This proves the
adaptive lower bounds and the mean-width lower bound.

The randomized batch policy omitting each $b_j$ with probability $1/K$
has expected width exactly $1/K$. Its expected $p$th power is
$K^{-1}\sum_j a_j^p\le1/K$, since $0\le a_j\le1$ and $p>1$.
For the reverse moment bound, average over the $K$ records
$a=e_j$, $b=0$. A batch omitting a $b_j$ has average $p$th-power width
$1/K$; one omitting an explicit $a_i$ has average $2/K$.
Therefore every batch distribution has worst-case moment at least $1/K$.

For precision, omitting a uniformly random $b_j$ succeeds with
probability at least $1/K$ at every record, because at least one capacity
is no greater than $1/K$. To prove the matching upper bound, use $K$
records indexed by $j$, with $b=0$,
\[
 a_j=\frac{1}{2K},\qquad
 a_i=\frac{2K-1}{2K(K-1)}\quad(i\ne j).
\]
These capacities sum to one. Only the batch omitting $b_j$ reaches
width at most $1/K$ at record $j$: every other $b$-omitting batch has larger
width, and every $a$-omitting batch has width $a_i+a_K>1/K$.
If $q_j$ is the probability of omitting $b_j$, worst-case success is
at most $\min_j q_j\le1/K$. The adaptive policy always succeeds.
The precision upper bound uses these interior capacity records;
checking only vertices would not establish it.
\end{proof}

The moment ratio is $K^{p-1}$, while the ratio of worst-case success
probabilities at precision $1/K$ is $K$. These separations use a
structured constructed feasible set and a budget querying all but one
coordinate. They establish the possible difference between objectives;
they do not establish comparable gains at small relative budgets in a
deployment or in the historical forecasting study.
